\documentclass[report]{itvest}

% Document-specific metadata
\title{Portfolio 2: Cloud Continuum}
\reportsubtitle{Case: Short Anonymised Case Name}
\author{Student Name} % group: \author{Name 1 \and Name 2 \and Name 3}
\date{\today}

% Portfolios use sections only, numbered 1, 2, ...
\renewcommand{\thesection}{\arabic{section}}

\begin{document}

\maketitle

\begin{checkit}
  \textbf{Assessment criteria} (delete this box before submitting):
  \begin{itemize}
    \item Placement decisions follow from the P1 quality attributes and
      constraints.
    \item Data flows respect the stated sovereignty rules, and the design
      enforces them.
    \item At least one security and one reliability tactic are
      demonstrated, not only described.
    \item A model can be rebuilt from Git, DVC and an MLflow run id, and
      one command deploys the system.
  \end{itemize}
\end{checkit}

\section{Deployment View}

Show which component runs on which tier (edge, on-premises, cloud), and
trace each placement to the P1 scenarios. Say how heterogeneous compute
(small devices, CPUs, GPUs) constrains the choice of model.

\begin{figure}[h]
  \centering
  % Replace the box with: \includegraphics[width=\textwidth]{deployment}
  \fbox{\parbox[c][7cm][c]{0.9\textwidth}{\centering
    Replace with the deployment view: components on tiers.}}
  \caption{Deployment view.}
  \label{fig:deployment}
\end{figure}

\begin{table}[h]
  \centering
  \begin{tabularx}{\textwidth}{llXl}
    \toprule
    Component & Tier & Why here & Scenario \\
    \midrule
    \ldots & \ldots & \ldots & QA-\ldots \\
    \ldots & \ldots & \ldots & QA-\ldots \\
    \bottomrule
  \end{tabularx}
  \caption{Component placement.}
  \label{tab:placement}
\end{table}

\section{Data Placement}

For each data item, state where it may be stored and processed, and the
rule behind it (contract, GDPR, customer policy). Say how the design
enforces the rule.

\begin{table}[h]
  \centering
  \begin{tabularx}{\textwidth}{XXXX}
    \toprule
    Data item & Stored at & Processed at & Rule \\
    \midrule
    \ldots & \ldots & \ldots & \ldots \\
    \ldots & \ldots & \ldots & \ldots \\
    \bottomrule
  \end{tabularx}
  \caption{Data placement.}
  \label{tab:data-placement}
\end{table}

\section{Build or Buy}

Compare at least one managed cloud or vendor ML platform with
self-hosting~\cite{bass_engineering_2025}, and record the decision as an
ADR.

\begin{table}[h]
  \centering
  \begin{tabularx}{\textwidth}{lXX}
    \toprule
    Criterion & Managed / vendor option & Self-hosted \\
    \midrule
    Cost & \ldots & \ldots \\
    Lock-in & \ldots & \ldots \\
    Skills & \ldots & \ldots \\
    Operations & \ldots & \ldots \\
    \bottomrule
  \end{tabularx}
  \caption{Managed versus self-hosted.}
  \label{tab:build-buy}
\end{table}

\subsection{ADR-3: Title of the Decision}

\begin{description}
  \item[Status:] Proposed / accepted.
  \item[Context:] The problem and the scenarios (QA-\ldots) that drive it.
  \item[Alternatives:] The managed or vendor option, and self-hosting.
  \item[Decision:] The chosen alternative.
  \item[Consequences:] The trade-off that decided it, and what gets harder.
\end{description}

\section{Security and Reliability Tactics}

Map each tactic to the component it protects. Mark which tactics are
demonstrated in Section~\ref{sec:evidence}. Say what happens when a part
fails or the connection drops.

\begin{table}[h]
  \centering
  \begin{tabularx}{\textwidth}{XlXll}
    \toprule
    Tactic & Type & Component & Scenario & Demonstrated \\
    \midrule
    \ldots & Security & \ldots & QA-\ldots & Yes / no \\
    \ldots & Reliability & \ldots & QA-\ldots & Yes / no \\
    \bottomrule
  \end{tabularx}
  \caption{Security and reliability tactics.}
  \label{tab:tactics}
\end{table}

\section{Deployment and Versioning}

Describe the Kubernetes manifests or Helm chart, the CI and deployment
workflow, and how data and models are versioned. Show how a model is
rebuilt from a Git commit, a DVC data version and an MLflow run id, and
the one command that deploys the system.

\section{Evidence}
\label{sec:evidence}

State whether the evidence comes from your own case or from the reference
project. Give the repository (a private repository or a zip upload is
fine), the output of \texttt{kubectl get all}, and a deleted pod or a
failed model showing recovery.

\begin{verbatim}
$ kubectl get all
...
\end{verbatim}

\section{AI Usage}

Briefly describe your use of AI tools on this assignment:

\begin{itemize}
  \item \emph{Code generation:} \ldots
  \item \emph{Grammar and writing:} \ldots
  \item \emph{Clarification of technical aspects:} \ldots
\end{itemize}

\bibliographystyle{IEEEtran}
\bibliography{bib}

\end{document}
